\documentclass{article}

\usepackage[preprint]{neurips_2025}

\usepackage[utf8]{inputenc}
\usepackage[T1]{fontenc}
\usepackage{hyperref}
\usepackage{url}
\usepackage{booktabs}
\usepackage{amsfonts}
\usepackage{amsmath}
\usepackage{nicefrac}
\usepackage{microtype}
\usepackage{xcolor}
\usepackage{enumitem}
\usepackage{graphicx}
\usepackage{subcaption}
\usepackage{algorithm}
\usepackage{algorithmic}
\usepackage{multirow}

\title{RL-based Dynamic Scheduling for LLM Serving:\\Improving Llumnix with Reinforcement Learning}

\author{
  Yuxuan Bai \\
  University of Virginia \\
  \texttt{yuxuan.bai@virginia.edu} \\
  \And
  Jiaqi Liu \\
  University of Virginia \\
  \texttt{jiaqi.liu@virginia.edu} \\
  \And
  Peng Xia \\
  University of Virginia \\
  \texttt{peng.xia@virginia.edu} \\
}

\begin{document}

\maketitle

\begin{abstract}
% TODO: Abstract (1 paragraph)
% Serving large language models at scale requires dynamic scheduling to handle
% unpredictable request lengths and KV cache growth. Llumnix (OSDI 2024) introduced
% live request migration with a heuristic virtual usage policy. We replace this
% heuristic with a reinforcement learning (RL) agent that learns scheduling decisions
% from data. Our PPO-based policy observes cluster state and outputs virtual usage
% values as a drop-in replacement for Llumnix's CalcVirtualUsage() function.
% Experiments on [X] GPUs with LLaMA-7B show that our RL policy achieves [X]%
% improvement in P99 latency and [X]% reduction in preemptions compared to the
% original heuristic, while generalizing across workload distributions.
\end{abstract}

%% ============================================================
\section{Introduction}
%% ============================================================

% TODO: 1 page
% Para 1: LLM serving is important, scheduling is hard (unpredictable lengths, KV cache, preemptions)
% Para 2: Llumnix solves this with migration + virtual usage heuristic
% Para 3: But heuristic is static, myopic, fixed tradeoffs
% Para 4: We propose RL-based policy. Key insight: minimal modification (only CalcVirtualUsage)
% Para 5: Results preview: X% improvement on Y metric
% Para 6: Contributions bullet list

%% ============================================================
\section{Background and Related Work}
%% ============================================================

% TODO: 1.5 pages

\subsection{LLM Serving and KV Cache Management}
% Auto-regressive inference, PagedAttention, preemption problem

\subsection{Llumnix: Dynamic Scheduling via Virtual Usage}
% Virtual usage abstraction, freeness metric, 4 rules
% Algorithm 1 from paper
% Migration mechanism

\subsection{Reinforcement Learning for Systems}
% Pensieve (video streaming), Decima (job scheduling), Resource management (Mao et al.)
% Why RL is well-suited: MDP structure, long-term optimization, adaptive

%% ============================================================
\section{Method}
%% ============================================================

% TODO: 2 pages

\subsection{MDP Formulation}

\paragraph{State space.}
% Mathematical definition of s_t
% Per-instance features, cluster features
% $s_t = \{s_t^{inst}, s_t^{cluster}\}$

\paragraph{Action space.}
% $a_t \in [-1, 1]^N$ where N = number of instances
% Maps to virtual usage: $v_i = p_i \cdot (1 + a_i \cdot \alpha)$

\paragraph{Reward function.}
% $r_t = -w_1 \cdot L_{P99} - w_2 \cdot N_{preempt} - w_3 \cdot F_{frag} + w_4 \cdot R_{SLO} + w_5 \cdot T$

\subsection{Policy Architecture}

% TODO: Include architecture figure
% \begin{figure}[t]
%   \centering
%   \includegraphics[width=0.9\textwidth]{figures/policy_architecture.pdf}
%   \caption{Multi-head policy network architecture.}
%   \label{fig:policy}
% \end{figure}

% Instance encoder → Self-attention → Cluster encoder → Combination → Action heads

\subsection{Training Procedure}

\paragraph{Phase 1: Behavior cloning.}
% Pre-train from heuristic traces

\paragraph{Phase 2: PPO fine-tuning.}
% Online training in simulated environment
% Hyperparameters table

% \begin{table}[h]
%   \centering
%   \caption{Training hyperparameters.}
%   \begin{tabular}{ll}
%     \toprule
%     Parameter & Value \\
%     \midrule
%     Algorithm & PPO \\
%     Learning rate & $3 \times 10^{-4}$ \\
%     Batch size & 64 \\
%     GAE $\lambda$ & 0.95 \\
%     Discount $\gamma$ & 0.99 \\
%     Clip range & 0.2 \\
%     Entropy coefficient & 0.01 \\
%     Total timesteps & $10^6$ \\
%     \bottomrule
%   \end{tabular}
% \end{table}

\subsection{Integration with Llumnix}
% RLBasedLoad class
% What changes vs. what stays the same
% Inference latency requirements (< 1ms)

%% ============================================================
\section{Experimental Setup}
%% ============================================================

% TODO: 1 page

\subsection{Hardware and Models}
% GPU type and count
% LLaMA-7B, 16-bit precision
% vLLM 0.6.3.post1

\subsection{Traces and Workloads}
% ShareGPT, BurstGPT
% Synthetic: S-S, M-M, L-L, S-L, L-S
% Arrival: Poisson, Gamma (varying CV)

% \begin{table}[h]
%   \centering
%   \caption{Sequence length distributions (number of tokens).}
%   \begin{tabular}{llccccc}
%     \toprule
%     & & Mean & P50 & P80 & P95 & P99 \\
%     \midrule
%     \multirow{2}{*}{ShareGPT} & In & 306 & 74 & 348 & 1484 & 3388 \\
%     & Out & 500 & 487 & 781 & 988 & 1234 \\
%     \midrule
%     \multirow{2}{*}{BurstGPT} & In & 830 & 582 & 1427 & 2345 & 3549 \\
%     & Out & 243 & 434 & 434 & 669 & 964 \\
%     \bottomrule
%   \end{tabular}
% \end{table}

\subsection{Baselines}
% Round-Robin, INFaaS++, Llumnix-Heuristic

\subsection{Metrics}
% End-to-end, prefill, decode latency (mean, P99)
% Preemption loss
% Fragmentation ratio

%% ============================================================
\section{Results}
%% ============================================================

% TODO: 2 pages

\subsection{Baseline Serving Performance (Exp 1)}
% ShareGPT + BurstGPT comparison
% Figure: 7-column latency plot

\subsection{Priority Workloads (Exp 2)}
% High-priority vs normal latencies
% SLO isolation

\subsection{Bursty Traffic (Exp 3)}
% Gamma arrivals with varying CV
% Adaptation capability

\subsection{Long-Context Workloads (Exp 4)}
% Memory fragmentation results
% L-L distribution

\subsection{Generalization (Exp 5)}
% Train on ShareGPT, test on BurstGPT
% Cross-distribution transfer

\subsection{Ablation Studies}
% Reward weight sensitivity
% State feature importance
% Training duration effect

%% ============================================================
\section{Discussion}
%% ============================================================

% TODO: 0.5 page
% When does RL help? When does it hurt?
% Training cost vs deployment benefit
% Inference overhead analysis
% Simulation vs. real deployment gap (if applicable)

%% ============================================================
\section{Limitations and Future Work}
%% ============================================================

% TODO: 0.5 page
% Simulation fidelity
% Scalability to larger clusters
% Multi-objective RL (Pareto optimization)
% Direct migration actions (Option B)
% Online adaptation

%% ============================================================
\section{Conclusion}
%% ============================================================

% TODO: 0.25 page
% Summary of contributions and key findings

%% ============================================================
\begin{thebibliography}{10}
\small

\bibitem[Sun et al.(2024)]{sun2024llumnix}
B.~Sun, Z.~Huang, H.~Zhao, W.~Xiao, X.~Zhang, Y.~Li, and W.~Lin.
Llumnix: Dynamic scheduling for large language model serving.
In \emph{OSDI}, 2024.

\bibitem[Kwon et al.(2023)]{kwon2023vllm}
W.~Kwon, Z.~Li, S.~Zhuang, Y.~Sheng, L.~Zheng, C.~H.~Yu, J.~E.~Gonzalez, H.~Zhang, and I.~Stoica.
Efficient memory management for large language model serving with {PagedAttention}.
In \emph{SOSP}, 2023.

\bibitem[Mao et al.(2016)]{mao2016resource}
H.~Mao, M.~Alizadeh, I.~Menache, and S.~Kandula.
Resource management with deep reinforcement learning.
In \emph{HotNets}, 2016.

\bibitem[Mao et al.(2017)]{mao2017neural}
H.~Mao, R.~Netravali, and M.~Alizadeh.
Neural adaptive video streaming with {Pensieve}.
In \emph{SIGCOMM}, 2017.

\bibitem[Schulman et al.(2017)]{schulman2017proximal}
J.~Schulman, F.~Wolski, P.~Dhariwal, A.~Radford, and O.~Klimov.
Proximal policy optimization algorithms.
\emph{arXiv:1707.06347}, 2017.

\bibitem[Mao et al.(2019)]{mao2019learning}
H.~Mao, M.~Schwarzkopf, S.~Venkatakrishnan, Z.~Meng, and M.~Alizadeh.
Learning scheduling algorithms for data processing clusters.
In \emph{SIGCOMM}, 2019.

% Add more references as needed

\end{thebibliography}

\end{document}
